\chapter{Gradient methods: from calculus to learning}
\label{ch:gradient-descent}

\begin{learningobjectives}
After this chapter, you should be able to:
\begin{itemize}
  \item derive gradient descent from a local linear approximation;
  \item explain directional derivatives, steepest descent, curvature, and the
  role of the learning rate;
  \item derive batch, stochastic, and mini-batch updates for least squares;
  \item connect a single identity-activation neuron to linear regression; and
  \item diagnose divergence, slow convergence, scaling problems, and false
  claims based only on optimizer convergence.
\end{itemize}
\end{learningobjectives}

\begin{chapterconnection}
Chapter~\ref{ch:linear-regression} described least squares as estimation and
projection. This chapter opens the numerical mechanism. The separation is
deliberate: a statistical model defines what is estimated; an optimizer defines
how parameters are searched. Gradient methods will return for logistic models,
matrix factorization, neural networks, inverse problems, and differentiable
scientific programs.
\end{chapterconnection}

\section{From steepest descent to stochastic approximation}

Iterative descent has roots in numerical analysis rather than in modern neural
network terminology. In 1847 Augustin-Louis Cauchy described a method of moving
along steepest descent while solving systems arising from quadratic
problems.\footnote{A.-L. Cauchy, ``M\'ethode g\'en\'erale pour la r\'esolution
des syst\`emes d'\'equations simultan\'ees,'' \emph{Comptes rendus de
l'Acad\'emie des Sciences}, 25, 1847.} In 1951 Herbert Robbins and Sutton Monro
gave a stochastic approximation procedure for finding a solution when only
noisy experimental responses were available.\footnote{H. Robbins and S.
Monro, ``A Stochastic Approximation Method,'' \emph{The Annals of Mathematical
Statistics}, 22(3), 1951,
\url{https://doi.org/10.1214/aoms/1177729586}.}

Modern gradient-based learning joins these two ideas: use local derivative
information to improve an objective, and estimate that information from a
manageable sample when the complete sum is expensive. Automatic
differentiation changes how derivatives are computed, not what optimization
question they answer.

\section{One variable: the derivative proposes a local move}

For differentiable $J:\mathbb{R}\rightarrow\mathbb{R}$, Taylor's formula gives
\[
J(w+\Delta w)=J(w)+J'(w)\Delta w+O((\Delta w)^2).
\]
When $J'(w)>0$, a small negative $\Delta w$ lowers the linear approximation;
when $J'(w)<0$, a small positive move does. Choosing
\[
\Delta w=-\eta J'(w),\qquad \eta>0,
\]
produces the gradient-descent iteration
\[
w_{t+1}=w_t-\eta_tJ'(w_t).
\]

The derivative supplies local direction and sensitivity, not a complete route
to a global minimizer. The learning rate $\eta_t$ converts sensitivity into a
step. If the step is too large, the neglected curvature term can dominate and
the objective may increase.

For the quadratic $J(w)=\tfrac12 a(w-w^*)^2$ with $a>0$,
\[
w_{t+1}-w^*=(1-\eta a)(w_t-w^*).
\]
The iterates converge exactly when $|1-\eta a|<1$, or
\[
0<\eta<\frac{2}{a}.
\]
This one-line calculation explains oscillation, divergence, and slow progress.
When $\eta$ is near $1/a$, progress is rapid; when it is tiny, each step retains
almost all of the preceding error.

\conceptcheck{For $J(w)=5(w-3)^2/2$, classify the behavior expected at
$\eta=0.05$, $0.2$, $0.35$, and $0.45$. Which choice reaches the minimizer in
one exact step?}

\section{Several variables: the gradient represents all directions}

For $J:\mathbb{R}^p\rightarrow\mathbb{R}$, the local approximation is
\[
J(\theta+d)\approx J(\theta)+\nabla J(\theta)^Td.
\]
For a unit direction $u$, the directional derivative is
\[
D_uJ(\theta)=\nabla J(\theta)^Tu.
\]
By Cauchy-Schwarz, this quantity is smallest over $\lVert u\rVert_2=1$ when
$u=-\nabla J/\lVert\nabla J\rVert_2$. Thus the negative gradient is the
direction of steepest local decrease under Euclidean distance.

The phrase ``steepest'' depends on geometry. Rescaling parameters changes
Euclidean distance and therefore changes the gradient path. A preconditioner,
natural gradient, or adaptive method corresponds to a different geometry or
coordinate-dependent scaling. This is why units and representation affect
optimization even when they leave the set of fitted predictions unchanged.

The Hessian
\[
H(\theta)=\nabla^2J(\theta)
\]
describes local curvature. Positive semidefinite $H$ throughout a convex domain
implies a convex objective. At a differentiable convex minimum the gradient is
zero, but a zero gradient in a nonconvex problem may also identify a saddle or
maximum.

\begin{misconceptionbox}
\textbf{Misconception:} the gradient points toward the nearest global minimum.
\textbf{Correction:} it points in the direction of steepest infinitesimal
decrease under a chosen norm. Curvature, constraints, step size, and nonconvex
geometry determine what happens after a finite step.
\end{misconceptionbox}

\section{Quadratic objectives make curvature visible}

Consider
\[
J(\theta)=\frac12\theta^TH\theta-c^T\theta+d,
\qquad H=H^T\succ 0.
\]
Then $\nabla J=H\theta-c$ and the unique minimizer is
$\theta^*=H^{-1}c$. Gradient descent satisfies
\[
\theta_{t+1}-\theta^*=(I-\eta H)(\theta_t-\theta^*).
\]
If $0<\lambda_{\min}\leq\lambda_j(H)\leq\lambda_{\max}$, a constant step
converges when
\[
0<\eta<\frac{2}{\lambda_{\max}}.
\]

The condition number $\kappa=\lambda_{\max}/\lambda_{\min}$ measures the
curvature disparity. A large $\kappa$ produces elongated contours. A safe step
for the steep direction makes little progress along the flat direction, so the
path zigzags slowly. Centering, scaling, reparameterization, and
preconditioning attempt to improve this geometry.

For an objective with $L$-Lipschitz gradient,
\[
J(v)\leq J(u)+\nabla J(u)^T(v-u)+\frac{L}{2}\lVert v-u\rVert_2^2.
\]
Setting $v=u-\eta\nabla J(u)$ shows that $0<\eta<2/L$ guarantees a decrease
unless the gradient is zero. This \term{descent lemma} generalizes the scalar
quadratic calculation.

\section{Least squares supplies an exact laboratory}

For design matrix $X\in\mathbb{R}^{n\times p}$ and target $y$, define
\[
J(w)=\frac{1}{2n}\lVert Xw-y\rVert_2^2.
\]
Differentiation gives
\[
\nabla J(w)=\frac1nX^T(Xw-y),
\qquad
\nabla^2J(w)=\frac1nX^TX.
\]
The Hessian is positive semidefinite, so every local minimum is global. If $X$
has full column rank, the minimum is unique. If not, gradient descent still
minimizes fitted loss but the coefficient selected depends on initialization
and the geometry of the updates.

The largest eigenvalue of $X^TX/n$ controls the stable constant step. Feature
scaling changes that spectrum. This gives a precise reason to scale before
gradient descent rather than the vague rule that models always require
standardization.

\begin{workedexample}{a linear neuron on real clinical measurements}
Let $x_i$ be a standardized baseline BMI measurement and let $y_i$ be a
one-year disease-progression score from the diabetes teaching table. A single
unit with identity activation computes
\[
z_i=wx_i+b,\qquad \widehat y_i=a(z_i)=z_i.
\]
With
\[
J(w,b)=\frac{1}{2n}\sum_{i=1}^n( wx_i+b-y_i)^2,
\]
the chain rule gives
\[
\frac{\partial J}{\partial w}
=\frac1n\sum_i(\widehat y_i-y_i)x_i,
\qquad
\frac{\partial J}{\partial b}
=\frac1n\sum_i(\widehat y_i-y_i).
\]
These are exactly the least-squares gradients. ``Single neuron'' and ``linear
regression'' describe two views of the same function and objective here. The
neuron language becomes informative when activation, composition, or shared
optimization later changes; it does not create a new estimator by itself.

Plot three synchronized objects: contours of $J(w,b)$ in parameter space, the
line $wx+b$ in data space, and loss against iteration. Parameter convergence,
visual fit, and training history answer different questions and should not be
collapsed into one plot.
\end{workedexample}

\section{Batch, stochastic, and mini-batch gradients}

If
\[
J(\theta)=\frac1n\sum_{i=1}^n\ell_i(\theta),
\]
batch gradient descent evaluates every $\nabla\ell_i$ before one update. A
stochastic update samples index $I_t$ and uses
\[
g_t=\nabla\ell_{I_t}(\theta_t).
\]
Under uniform sampling,
$\mathbb{E}[g_t\mid\theta_t]=\nabla J(\theta_t)$: the stochastic gradient is
unbiased, though noisy. A mini-batch averages several sampled gradients to
trade variance for parallel efficiency.

One pass through the training observations is an \term{epoch}. An epoch is not
one update unless the batch is the entire dataset. Shuffling affects the
sequence of stochastic gradients, so the seed and data order belong to the
reproducibility record.

Constant-step stochastic methods generally fluctuate near a solution because
gradient noise remains. Decreasing schedules can support convergence under
conditions; modern practice also uses warmup, momentum, adaptive scaling, and
averaging. Their convenience does not remove the need to monitor the actual
objective and update magnitude.

\begin{tabularx}{\textwidth}{@{}p{0.18\textwidth}p{0.25\textwidth}X@{}}
\toprule
Update & Main advantage & Main cost or risk \\
\midrule
Full batch & Exact gradient of the finite-sample objective & Expensive update;
limited streaming behavior \\
Stochastic & Cheap updates and immediate online use & High variance and strong
sensitivity to order and schedule \\
Mini-batch & Efficient vectorization with reduced variance & Batch size changes
optimization noise, memory use, and sometimes generalization behavior \\
\bottomrule
\end{tabularx}

\section{Applications beyond fitting a line}

Gradient methods appear whenever a useful objective is differentiable or has a
tractable derivative surrogate:

\begin{itemize}
  \item logistic regression minimizes a convex likelihood-based loss;
  \item matrix factorization estimates latent structure in recommendation and
  imaging;
  \item inverse problems adjust hidden physical parameters so a simulator
  matches measurements;
  \item optimal control and differentiable simulation optimize trajectories or
  designs;
  \item neural networks compose many parameterized maps and obtain derivatives
  by reverse-mode automatic differentiation; and
  \item variational scientific models optimize an energy or evidence bound.
\end{itemize}

Each application changes the meaning of the objective and the consequences of
local minima, constraints, and stochasticity. The update formula alone does not
validate the model.

\section{Stopping and diagnosing responsibly}

Useful diagnostics include training and validation objectives, gradient norm,
step norm, parameter norm, nonfinite values, and elapsed resources. A stopping
rule might require a small relative objective change and small gradient norm,
with a maximum iteration budget. Validation-based early stopping is also model
selection and must remain inside the development boundary.

\begin{warningbox}
Optimizer convergence establishes only that the implemented search reached a
stationary or sufficiently stable point for its objective. It does not show
that gradients are correct, the objective is scientifically appropriate, the
data are representative, or predictions generalize.
\end{warningbox}

For a hand-built optimizer, verify derivatives three ways when practical:

\begin{enumerate}
  \item derive them symbolically with shapes shown;
  \item compare against centered finite differences on a tiny problem; and
  \item compare the final fitted result with a trusted solver on a convex case.
\end{enumerate}

Finite differences are diagnostic, not the preferred large-scale training
method. Their truncation and floating-point errors depend on perturbation size.

\begin{professionalbox}
A reusable optimizer should report termination reason, iterations, final
objective, gradient or update norm, and nonfinite failure. Record the objective,
data version, preprocessing, initialization, seed, batch construction,
learning-rate schedule, precision, and software version. Silently returning the
last iterate after divergence is an interface failure.
\end{professionalbox}

\section{Exercises}

\begin{enumerate}
  \item Derive the exact scalar quadratic convergence interval and describe the
  paths for positive and negative contraction factors.
  \item Prove that $-\nabla J/\lVert\nabla J\rVert_2$ minimizes the directional
  derivative over Euclidean unit vectors.
  \item For a two-feature least-squares problem, compute $X^TX/n$, its
  eigenvalues, condition number, and stable constant-step interval before and
  after standardization.
  \item Derive the weight and intercept gradients for a mini-batch linear
  neuron, with every array shape annotated.
  \item Construct a case where training loss converges while validation loss
  worsens. Explain why changing the stopping rule is model selection.
  \item Design tests that catch an incorrect sign, missing factor, silent
  nonfinite update, and mutation of the input arrays.
\end{enumerate}

\begin{chapterreview}
Gradient descent follows the negative gradient because it minimizes the local
linear approximation under Euclidean geometry. Curvature determines whether a
finite step decreases the objective and how rapidly different directions
converge. Least squares exposes this behavior exactly through $X^TX/n$ and
connects an identity-activation neuron to a familiar estimator. Stochastic and
mini-batch methods estimate the full gradient more cheaply, but introduce
variance and additional reproducibility choices. Optimization success remains
distinct from statistical generalization and scientific validity.
\end{chapterreview}

\begin{notebooklink}
Lecture 10 notebook 00 begins with tangent lines and hand-computed scalar
updates, moves to contours and directional derivatives, then trains a single
identity-activation neuron on real diabetes observations. Synchronized plots
connect parameter space, data space, and iteration space; finite differences
and a trusted least-squares solver verify the implementation. Scaling is then
explained through Hessian eigenvalues and condition numbers.
\end{notebooklink}
